\begin{pastexam}[label={px:ncku-106-cce-ds-06}]{成大電通所 106 學年度 第二大題第 2 題}
（簡答；15 分）

\textbf{(15 pts) Transitive closure graph.}

According to the graph G shown in the following:

\begin{center}
\begin{tikzpicture}[v/.style={draw,circle,minimum size=6mm,inner sep=0pt,font=\small},>={Stealth[length=2.4mm]},thick]
  \node[v] (1) at (0,0) {1};
  \node[v] (2) at (1.6,1.0) {2};
  \node[v] (3) at (1.6,-1.0) {3};
  \node[v] (4) at (3.8,1.0) {4};
  \node[v] (5) at (3.8,-1.0) {5};
  \draw[->] (1) -- (2);
  \draw[->] (1) -- (3);
  \draw[->] (2) -- (4);
  \draw[->] (3) -- (5);
  \draw[->] (4) -- (3);
  \draw[->] (5) -- (2);
\end{tikzpicture}
\end{center}
\begin{enumerate}[label=\alph*.]
\item \textbf{(1 pts)} Use an adjacency matrix to represent the graph.
\item \textbf{(6 pts)} Find the transitive closure matrix A$^+$ and the reflexive transitive closure
matrix A$^*$ of G.
\item \textbf{(2 pts)} Explain the meaning of A$^+$.
\item \textbf{(2 pts)} Explain the meaning of A$^*$.
\item \textbf{(4 pts)} If we use a matrix A$^0$ to represent the distances between any two
vertices, where t A$^0$(i, j) denotes the shortest distance between two vertices i and j,
Floyd-Warshall algorithm can find the shortest path between any two vertices by compute a new
matrix A$^k$ from A$^{k-1}$. Please list the general recursive equation and explain the equation
shortly (hint: what does k represent in the equation?).
\end{enumerate}
\end{pastexam}

\begin{pastsoln}
圖的邊：$1 \to 2$、$1 \to 3$、$2 \to 4$、$3 \to 5$、$4 \to 3$、$5 \to 2$。 其中
$2 \to 4 \to 3 \to 5 \to 2$ 是一個環。

\textbf{a.} 鄰接矩陣 $A$（$A[i][j] = 1$ 表示有邊 $i \to j$，資結書第~21.2~節）：
\[
A = \begin{pmatrix}
0 & 1 & 1 & 0 & 0 \\
0 & 0 & 0 & 1 & 0 \\
0 & 0 & 0 & 0 & 1 \\
0 & 0 & 1 & 0 & 0 \\
0 & 1 & 0 & 0 & 0
\end{pmatrix}
\]

\textbf{b.} 逐一看每個頂點走得到哪些頂點（至少走一條邊）：
\begin{itemize}
\item $1$：$1 \to 2$、$1 \to 3$，再沿著環走到 $4$、$5$；沒有邊進入 $1$，所以回不到 $1$。 可到 $\{2, 3, 4, 5\}$。
\item $2, 3, 4, 5$：都在環 $2 \to 4 \to 3 \to 5 \to 2$ 上，沿著環可以走到環上每一個頂點（包括
繞一圈回到自己），但走不到 $1$。 都可到 $\{2, 3, 4, 5\}$。
\end{itemize}
\[
A^{+} = \begin{pmatrix}
0 & 1 & 1 & 1 & 1 \\
0 & 1 & 1 & 1 & 1 \\
0 & 1 & 1 & 1 & 1 \\
0 & 1 & 1 & 1 & 1 \\
0 & 1 & 1 & 1 & 1
\end{pmatrix},
\qquad
A^{*} = \begin{pmatrix}
1 & 1 & 1 & 1 & 1 \\
0 & 1 & 1 & 1 & 1 \\
0 & 1 & 1 & 1 & 1 \\
0 & 1 & 1 & 1 & 1 \\
0 & 1 & 1 & 1 & 1
\end{pmatrix}
\]
$A^{*}$ 是 $A^{+}$ 再把對角線全設成 $1$；只有 $A^{*}[1][1]$ 和 $A^{+}$ 不同，因為 $2$--$5$ 在環上，
$A^{+}$ 的對角線本來就是 $1$。 也可以照資結書虛擬碼~25.10（Warshall）計算：
$A^{(k)}[i][j] = A^{(k-1)}[i][j] \lor (A^{(k-1)}[i][k] \land A^{(k-1)}[k][j])$，$k = 1, \ldots, 5$，
結果相同。

\textbf{c.} $A^{+}[i][j] = 1$ 若且唯若從 $i$ 到 $j$ 有一條長度\textbf{至少為 $1$} 的路徑。 所以
$A^{+}[i][i] = 1$ 表示 $i$ 在某個環上。

\textbf{d.} $A^{*}[i][j] = 1$ 若且唯若從 $i$ 到 $j$ 有一條長度\textbf{至少為 $0$} 的路徑，也就是
$i = j$ 或 $A^{+}[i][j] = 1$。 每個頂點都算走得到自己（自反），對角線全是 $1$。

\textbf{e.} 題目說「$A^{0}(i, j)$ 表示 $i$、$j$ 之間的最短距離」不太精確：$A^{0}$ 是起始的距離矩陣，
$A^{0}(i, j)$ 是邊 $i \to j$ 的長度（沒有邊時是 $\infty$，$A^{0}(i, i) = 0$），最後的 $A^{n}$ 才是
最短距離。 遞迴式（資結書第~25.6~節）：
\[
A^{k}(i, j) = \min\bigl\{\, A^{k-1}(i, j),\; A^{k-1}(i, k) + A^{k-1}(k, j) \,\bigr\},
\qquad k = 1, 2, \ldots, n .
\]
$A^{k}(i, j)$ 是從 $i$ 到 $j$、中間只准經過編號 $1, 2, \ldots, k$ 的頂點時的最短距離。 從 $A^{k-1}$
到 $A^{k}$ 時，多允許頂點 $k$ 當中間點：最短路徑要嘛不經過 $k$，長度是 $A^{k-1}(i, j)$；要嘛經過
$k$ 恰好一次（沒有負環時），拆成 $i \leadsto k$ 和 $k \leadsto j$ 兩段，兩段的中間點都只在
$\{1, \ldots, k - 1\}$ 裡，長度是 $A^{k-1}(i, k) + A^{k-1}(k, j)$。 兩者取小。 $k$ 從 $1$ 做到 $n$
後，所有頂點都可以當中間點，$A^{n}(i, j)$ 就是最短距離，總時間 $\Theta(n^3)$。
\end{pastsoln}
